\section{Conventions and preliminary notation}\label{sec:conventions}

\subsection{Pairs, nef data and discrepancies}

All varieties are integral. A contraction is a projective surjective
morphism $f$ to a normal variety with $f_*\cO=\cO$. Unless a different
field is specified, the local and field-theoretic arguments take place
over $\C$. Section~\ref{sec:global} explains the passage to any
algebraically closed field of characteristic zero. Auxiliary function
fields are finitely generated over $\C$ and need not be algebraically
closed.

For a generalized pair $(X,B+M)$, choose a projective birational model
$\pi\colon X'\to X$ carrying the nef rational Cartier divisor $M'$.
On a smooth model $p\colon W\to X$ dominating $X'$, write $M_W$ for
its pullback and define $B_W$ by
\[
 K_W+B_W+M_W=p^*(K_X+B+M).
\]
For a prime divisor $E$ on $W$, set
\[
 A_{X,B+M}(E)=1-\operatorname{coeff}_E(B_W).
\]
This is independent of the choice of higher model. The nef data are
held fixed throughout all birational operations on a generalized pair.
An ordinary subpair may have negative boundary coefficients; its
discrepancy is defined by the same crepant formula without a nef part.
The terms lc and klt refer respectively to nonnegative and positive
log discrepancies. Effectivity, when needed, is stated separately.

\subsection{Homogeneous valuations and volume forms}

A divisorial valuation means a positive real multiple
$v=c\ord_E$ of a prime-divisor order, unless an integral or rational
normalization is specified. Discrepancies and orders of divisors are
extended homogeneously. An integral-valued divisorial valuation need
not be primitive. We also use quasi-monomial valuations on fixed
log-smooth models, where positive weights on boundary parameters define
their values; at a face, zero-weight directions belong to the residue
coefficient field. Any finite-group action on functions and
differentials uses the same pullback convention. The angle of a
character takes values in $\R/\Z$, with exponential
$a\mapsto \exp(2\pi i a)$.

For a real valuation $v$ centred at a closed point $x$ of a normal
$n$-fold, $n\ge1$, put
\[
 \mathfrak a_t(v)=\{g\in\cO_{X,x}:v(g)\ge t\},\qquad
 \mathfrak a_{>t}(v)=\{g\in\cO_{X,x}:v(g)>t\}.
\]
These ideals have finite colength for $t>0$: sufficiently high powers
of the maximal ideal lie in them, since its finitely many generators
have positive value. We use
\[
 \vol_{X,x}(v)=\limsup_{t\to\infty}
 \frac{n!\,\operatorname{length}(\cO_{X,x}/\mathfrak a_t(v))}{t^n},
 \qquad
 \gr_v\cO_{X,x}=\bigoplus_{t\ge0}
             \mathfrak a_t(v)/\mathfrak a_{>t}(v).
\]
For a klt pair $(X,\Delta)$ and finite log discrepancy, its normalized
volume is $\widehat{\vol}_{X,\Delta}(v)=A_{X,\Delta}(v)^n\vol_{X,x}(v)$;
it is $+\infty$ when the discrepancy is infinite
\cite[Definitions~2.4--2.5]{MinimizerUniqueness}.
Thus $\vol(cv)=c^{-n}\vol(v)$ for $c>0$, and normalized volume is
invariant under positive rescaling. The notation $(v>T)$ means
$\mathfrak a_{>T}(v)$. The inclusions
$\mathfrak a_{t+1}\subset\mathfrak a_{>t}\subset\mathfrak a_t$
show that strict cutoffs give the same normalized limsup. For the
finitely generated positive rational gradings used below, Hilbert
asymptotics give the limit and the leading colength coefficient
$\vol(v)/n!$. This numerical valuation volume is distinct from a
differential volume form and from the volume of a divisor.

Let $L/\C$ be a function field of transcendence degree $d$. An
$r$-pluri-volume form is a nonzero element of
$(\bigwedge^d\Omega^1_{L/\C})^{\otimes r}$, with $r$ a positive
integer. On a smooth model its associated subpair has boundary
$-r^{-1}\operatorname{div}(\omega)$; write $A_\omega$ for the
resulting homogeneous discrepancy. For $r=1$ we say volume form.
For a volume form $\eta$, the index-one formulas are
\begin{equation}\label{glob:form-discrepancy}
 A_\eta(c\ord_E)=c\bigl(1+\ord_E(\eta)\bigr),
 \qquad A_{h\eta}(u)=A_\eta(u)+u(h).
\end{equation}
If $v$ is integral-valued, the discrepancy of a volume form is
integral. This index-one fact is used in the integral jump argument.

If two forms $\omega,\sigma$ are compared at a common positive index
$p$ and $H=\omega^{\otimes p/r_\omega}/
\sigma^{\otimes p/r_\sigma}$, then
\[
 A_\omega(v)-A_\sigma(v)=\frac{1}{p}v(H).
\]
Here the quotient denotes a rational function after identifying the
one-dimensional spaces of top forms. Indeed, the boundary difference
is $-p^{-1}\operatorname{div}(H)$, which gives the displayed sign.
Under a finite characteristic-zero field extension, the discrepancy
of a pulled-back form at an upstairs divisorial valuation equals that
of the restricted valuation downstairs. Indeed, if a primitive
valuation $\ord_{E'}$ upstairs restricts to $e\ord_E$ downstairs,
pullback of a volume form satisfies
\[
 \ord_{E'}(\eta)=e\ord_E(\eta)+(e-1).
\]
Adding $1$ gives $A_\eta(\ord_{E'})=eA_\eta(\ord_E)$; homogeneity
and tensor powers give the assertion for every normalization and
pluri-index. This identity already accounts for the scale of the
restricted valuation; no separate rescaling is implicit. The
restricted nontrivial valuation is divisorial, since the value-group
index and residue-field extension are finite.

For a subfield $K\subset L$ and a form on $L$, a barred discrepancy
denotes the infimum over divisorial prolongations with the specified
restriction to $K$, whenever introduced in the relevant setup. The
precise admissible valuations and hypotheses are stated in
Section~\ref{sec:integral-jump}. The notation $A_\omega$ always refers
to the form in that local setup, not to a fixed boundary on every model.

\subsection{Positivity, bounds and limits}

Divisor pullbacks are taken for rational Cartier multiples. A degree
on an embedded projective variety means intersection with the
appropriate power of its hyperplane class. When volumes of divisors
are used, roundings and divisibility are taken on the fixed model
specified before the limit. Constants in a sequential assertion may
become uniform after the stated passage to a subsequence. Each such
assertion specifies which dimension, bounded family, fixed comparison
constant or earlier choice is permitted to enter the bound.

An NQC nef part is a nonnegative real combination of nef rational
Cartier b-divisors. The rational nef data in
Theorem~\ref{thm:main} are therefore NQC. We use the standard klt
MMP, relative vanishing and big-boundary existence results as
established literature. All stronger uniformity, specialization,
valuation-selection and termination statements needed for the proof
are established in the sections that follow.
